% 第8节：相互作用Wightman函数
\section{相互作用Wightman函数}
\label{sec:wightman_functions}

\subsection{引言}

自由场Wightman函数正则量子化得到，而相互作用场需要微扰或非微扰构造。我们给出CTUFT的微扰框架。

\subsection{微扰论}

相互作用哈密顿量为：
\begin{equation}
    \hat{H}_{\text{int}} = \kappa_T \int d^3x \, \hat{\mathcal{T}}^3(x)
\end{equation}

\subsection{时间编序关联函数}

在相互作用绘景中：
\begin{equation}
    G_n(x_1, \ldots, x_n) = \langle 0 | T\{\hat{\mathcal{T}}_I(x_1) \cdots \hat{\mathcal{T}}_I(x_n) S\} | 0 \rangle
\end{equation}
其中$S = T\{\exp(-i\int dt \, \hat{H}_{\text{int}}(t))\}$。

\subsection{费曼规则}

\begin{enumerate}
    \item \textbf{传播子}：$i/(p^2 - M_T^2 + i\epsilon)$
    \item \textbf{顶点}：$-i\kappa_T (2\pi)^4 \delta^{(4)}(\sum p_i)$
    \item \textbf{外线}：1（截肢的）或$e^{-ip\cdot x}$（未截肢的）
\end{enumerate}

\subsection{进展状态（70\%）}

\subsubsection{已完成}
\begin{itemize}
    \item 两点函数到两圈阶
    \item 三点函数在树图阶
    \item 四点散射振幅
\end{itemize}

\subsubsection{进行中}
\begin{itemize}
    \item 高圈修正
    \item 非微扰效应（瞬子）
    \item 束缚态贡献
\end{itemize}

\subsection{谱性质}

相互作用两点函数的Källén-Lehmann表示：
\begin{equation}
    G_2(p) = \frac{Z}{p^2 - M_{\text{phys}}^2 + i\epsilon} + \int_{4M_T^2}^\infty d\mu^2 \frac{\rho(\mu^2)}{p^2 - \mu^2 + i\epsilon}
\end{equation}
其中$Z$是波函数重整化，$M_{\text{phys}}$是物理质量。
